\section*{Physique (13 points)}
\section*{Exercice 1 (3 points) : Radioactivité et médecine nucléaire}
La médecine nucléaire est la spécialité médicale qui utilise les propriétés de la radioactivité à des fins médicales. Un radionucléide quelconque a autant de chance de se désintégrer à un moment donné qu'un autre radionucléide de la même espèce, et la désintégration ne dépend pas des conditions physico-chimiques dans lesquelles le nucléide se trouve.

On dispose de trois échantillons contenant chacun un des radionucléides : le Samarium 153, l'Yttrium 90 et le Lutétium 177. Le nombre initial de noyaux dans chaque échantillon est le même:\\
$N_{0}\left({ }_{62}^{153} S m\right)=N_{0}\left({ }_{39}^{90} Y\right)=N_{0}\left({ }_{71}^{177} L u\right)$.\\
Les demi-vies des trois radionucléides sont telles que : $t_{1 / 2}\left({ }_{62}^{153} S m\right)<t_{1 / 2}\left({ }_{39}^{90} Y\right)<t_{1 / 2}\left({ }_{71}^{177} L u\right)$.

\section*{Données :}
\begin{center}
\begin{tabular}{|l|c|c|c|}
\cline { 2 - 4 }
\multicolumn{1}{c|}{} & 39 protons & 51 neutrons & Noyau ${ }_{39}^{90} Y$ \\
\hline
Énergie de masse en $(M e V)$ & 36592,8 & 47918,1 & 83748,5 \\
\hline
\end{tabular}
\end{center}

\begin{center}
\begin{tabular}{|c|c|c|}
\hline
Noyau & ${ }_{62}^{153} \mathrm{Sm}$ & ${ }_{71}^{177} \mathrm{Lu}$ \\
\hline
Énergie de liaison par nucléon en $(\mathrm{MeV} / \mathrm{nucléon})$ & 8,23 & 7,85 \\
\hline
\end{tabular}
\end{center}

\begin{enumerate}
  \item Définir la demi-vie d'un radionucléide.
  \item Le graphe ci-contre donne les courbes de décroissance radioactive des trois radionucléides cités.\\
2.1. Attribuer, en justifiant votre réponse, chaque courbe au radionucléide correspondant.\\
2.2. Déterminer graphiquement :
\end{enumerate}

\begin{itemize}
  \item le nombre initial $N_{0}\left({ }_{39}^{90} Y\right)$;
  \item la demi-vie du radionucléide ${ }_{39}^{90} Y$.\\
\includegraphics[max width=\textwidth, alt={}, center]{89340c76-9196-4142-968f-a7c291f2aa8d-4_646_1011_1484_959}
\end{itemize}

\begin{enumerate}
  \setcounter{enumi}{2}
  \item On désire comparer la stabilité des trois radionucléides.\\
3.1. Calculer, en unité $M e V$, l'énergie de liaison du radionucléide ${ }_{39}^{90} Y$.\\
3.2. Déterminer, en justifiant votre réponse, le noyau le plus stable parmi ${ }_{62}^{153} S m$, ${ }_{39}^{90} Y$ et ${ }_{71}^{177} \mathrm{Lu}$.
  \item Le radionucléide ${ }_{39}^{90} \mathrm{Y}$ se désintègre en Zirconium ${ }_{40}^{90} \mathrm{Zr}$.\\
4.1. Écrire l'équation de désintégration du radionucléide ${ }_{39}^{90} Y$.\\
4.2. Pour irradier à midi ( 12 h ) une tumeur hépatique, un médecin prescrit l'injection d'une dose d'Yttrium 90 d'activité $a=5 \cdot 10^{9} B q$. Cette dose a été préparée au laboratoire à 7h le même jour. Déterminer l'activité de l'Yttrium 90 à 7 h .
\end{enumerate}

\section*{Exercice 2 (5 points): Dipôle RL - Circuit RLC série}
La présence d'une bobine dans un circuit électrique alimenté par un générateur impose un comportement de celle-ci qui se traduit par une variation de l'intensité du courant. Lorsque la bobine est associée à un condensateur chargé et un conducteur ohmique, ces éléments peuvent constituer un oscillateur libre siège d'un échange énergétique et d'oscillations électriques qui peuvent être entretenues.

\section*{Cet exercice vise :}
\begin{itemize}
  \item l'étude de la réponse d'un dipôle RL soumis à un échelon de tension;
  \item l'étude énergétique d'un circuit RLC série.
\end{itemize}

\section*{Partie 1 : Étude d'un dipôle RL}
Pour étudier la réponse d'un dipôle $R L$ à un échelon de tension ascendant, on dispose du matériel suivant :

\begin{itemize}
  \item un générateur idéal de tension de force électromotrice $E=6 V$;
  \item un conducteur ohmique de résistance $R=10 \Omega$;
  \item une bobine d'inductance $L=10 m H$ et de résistance négligeable ;
  \item un interrupteur $K$.
\end{itemize}

\begin{enumerate}
  \item Proposer le schéma du montage expérimental permettant de réaliser cette étude.
  \item On ferme l'interrupteur $K$ à l'instant $t_{0}=0$. On note $i$ l'intensité du courant qui traverse le circuit. Représenter sur le schéma proposé la tension $u_{R}$ aux bornes du conducteur ohmique et la tension $u_{L}$ aux bornes de la bobine en convention récepteur.
  \item L'équation différentielle vérifiée par l'intensité $i$ s'écrit $\frac{d i}{d t}+\frac{1}{\tau} \cdot i=\frac{E}{A}$.\\
3.1. Déterminer les expressions des constantes $\tau$ et $A$.\\
3.2. En utilisant les équations aux dimensions, déterminer la dimension de $\tau$ et calculer sa valeur.
  \item La solution de l'équation différentielle s'écrit $i(t)=\frac{E}{R} \cdot\left(1-e^{-\frac{t}{\tau}}\right)$.\\
4.1. Recopier, sur votre copie, le numéro de la question et écrire la lettre correspondante à la proposition vraie.\\
L'expression de la tension aux bornes de la bobine en volt est :
\end{enumerate}

\begin{center}
\begin{tabular}{|l|l|l|l|l|l|l|l|}
\hline
$\mathbf{A}$ & $u_{L}(t)=6 .\left(1-e^{-10^{3} \cdot t}\right)$ & $\mathbf{B}$ & $u_{L}(t)=6 . e^{-10^{3} \cdot t}$ & $\mathbf{C}$ & $u_{L}(t)=0,6 . e^{-10^{-3} \cdot t}$ & $\mathbf{D}$ & $u_{L}(t)=6 . e^{-10^{-3} \cdot t}$ \\
\hline
\end{tabular}
\end{center}

4.2. Déterminer, en régime permanent, la valeur de l'intensité $I_{0}$ du courant électrique.\\
5. Quel rôle a joué la bobine durant la phase $0<t<5 . \tau$ ?

\section*{Partie 2: Étude énergétique d'un circuit RLC série}
On monte la bobine et le conducteur ohmique précédents en série avec un condensateur de capacité $C=1 \mu \mathrm{~F}$, initialement chargé . On ferme l'interrupteur $K$ à l'instant $t_{0}=0$.

\begin{enumerate}
  \item Établir l'équation différentielle vérifiée par la tension $u_{C}$ aux bornes du condensateur.
  \item Une étude expérimentale a permis de tracer les courbes de l'énergie électrique $\mathscr{E}_{e}$ emmagasinée dans le condensateur et de l'énergie totale $\mathscr{E}$ du circuit (figure ci-contre).\\
\includegraphics[max width=\textwidth, alt={}, center]{89340c76-9196-4142-968f-a7c291f2aa8d-5_645_959_2156_1027}
\end{enumerate}

En exploitant les courbes :\\
a. Déterminer la valeur de l'énergie totale $\mathscr{E}_{0}$ du circuit à l'instant $t_{0}=0$.

Déduire la valeur de la charge initiale $Q_{0}$ du condensateur à l'instant $t_{0}=0$.\\
b. Déterminer, à l'instant $t_{1}=0,9 \mathrm{~ms}$, la valeur de l'énergie électrique $\mathscr{E}_{e 1}$ emmagasinée dans le condensateur et la valeur de l'énergie totale $\mathscr{E}$ du circuit.\\
c. Déterminer la valeur absolue de l'intensité $i_{1}$ du courant électrique dans le circuit à l'instant $t_{1}$.\\
d. Expliquer la diminution de l'énergie totale du circuit.\\
3. Pour entretenir les oscillations électriques dans le circuit, on ajoute à celui-ci un générateur délivrant une tension $u_{G}(t)=k . i(t)$ avec $k$ constante positive.\\
3.1. Quelle doit être la valeur de $k$ ?\\
3.2. Calculer la valeur de la période des oscillations électriques dans ce cas.

\section*{Exercice 3 (5 points) : Parc de jeux}
Le parc de jeu constitue un cadre où se pratique du sport pour enfants. Les mouvements des enfants sont généralement de deux types et diffèrent par la nature des actions mécaniques exercées sur ces enfants. L'application des lois de Newton permet de déterminer l'évolution de certaines grandeurs cinématiques et dynamiques caractérisant les mouvements des enfants.

Cet exercice vise l'étude de deux types de mouvement et la détermination de certaines grandeurs qui les caractérisent.

Dans un parc de jeux, se trouve une glissière dont le profil est représenté dans le plan vertical (figure 1). La glissière est constituée :

\begin{itemize}
  \item d'une partie $A B$ rectiligne de longueur $A B=L$ inclinée d'un angle $\alpha$ par rapport à l'horizontal ;
  \item d'une partie de chute constituée d'un filet horizontal, situé à une hauteur $h$ de $B$.
\end{itemize}

On se propose d'étudier le mouvement du centre d'inertie $G$ d'un enfant de masse $m$ sur la glissière.

\begin{figure}[h]
\begin{center}
  \includegraphics[alt={},max width=\textwidth]{89340c76-9196-4142-968f-a7c291f2aa8d-6_508_716_1720_776}
\captionsetup{labelformat=empty}
\caption{Figure 1}
\end{center}
\end{figure}

\section*{1. Étude du mouvement sur $A B$}
L'enfant part de la position $A$ sans vitesse initiale. Les frottements sont équivalents à une force constante $\vec{f}$ de même direction que le vecteur vitesse et de sens opposé.\\
Pour étudier le mouvement de $G$, on choisit un repère ( $A, \vec{i}, \vec{j}$ ) lié à la Terre supposé galiléen, et l'instant de départ de $G$ de A comme origine de temps ( $t_{0}=0$ ). On repère la position de $G$ à un instant $t$ par son abscisse $X_{G}$ dans ce repère.\\
À $t_{0}=0, x_{G}=x_{0}=0$ (figure 1).

Données : $f=41 \mathrm{~N} ; m=20 \mathrm{~kg} ; g=10 \mathrm{~m} . \mathrm{s}^{-2} ; \alpha=40^{\circ}$\\
1.1. En appliquant la deuxième loi de Newton, établir l'équation différentielle vérifiée par l'abscisse $X_{G}$.\\
1.2. Déterminer, en justifiant, la nature du mouvement de $G$.\\
1.3. Lors du mouvement, $G$ passe par la position $B$ à l'instant $t_{B}=1,35 s$.\\
a. Calculer la distance $A B$.\\
b. Vérifier que la valeur de la vitesse en $B$ est $v_{B}=5,9 m \cdot s^{-1}$.\\
1.4. Déterminer l'intensité de la force $\vec{R}$ exercée par le plan incliné sur l'enfant.

\section*{2. Étude du mouvement de chute}
L'enfant quitte la partie $A B$ en $B$ avec la vitesse $\vec{v}_{B}$ et tombe sur le filet de protection.\\
Pour étudier le mouvement de chute, on choisit un autre repère ( $B, \vec{i}^{\prime}, \vec{j}^{\prime}$ ) lié à la Terre supposé galiléen, et l'instant de passage de $G$ en $B$ comme nouvelle origine de temps ( $t_{0}=0$ ).\\
2.1. Établir l'expression de l'équation de la trajectoire du mouvement de $G$ dans le $\operatorname{repère}\left(B, \vec{i}^{\prime}, \vec{j}^{\prime}\right)$. Déduire sa nature.\\
2.2. L'enfant tombe à l'instant $t_{P}$ sur le filet en une position où les coordonnées de $G$ sont : $\left(x_{P}=1,6 m \quad ; \quad y_{P}=2 m\right)$.\\
a. Calculer la valeur de $t_{P}$.\\
b. Déterminer la valeur de la vitesse de $G$ à l'instant $t_{P}$.\\
2.3. Un système d'acquisition convenable a permis de représenter les trajectoires (1), (2) et (3) des centres d'inertie de trois enfants $\left(E_{1}\right),\left(E_{2}\right)$ et $\left(E_{3}\right)$ qui ont quitté la position $B$ avec les vitesses respectives $v_{1}=3,5 \mathrm{~m} . \mathrm{s}^{-1}, v_{2}=4,5 \mathrm{~m} . \mathrm{s}^{-1}$ et $v_{3}=5,8 m \cdot s^{-1}$ (figure 2).\\
a. Attribuer à chaque enfant la trajectoire qui lui correspond.\\
b. Quel est l'enfant qui a eu la plus longue durée de chute ? Justifier.

\begin{figure}[h]
\begin{center}
  \includegraphics[alt={},max width=\textwidth]{89340c76-9196-4142-968f-a7c291f2aa8d-7_405_734_1457_1192}
\captionsetup{labelformat=empty}
\caption{Figure 2}
\end{center}
\end{figure}



